\documentclass[11pt]{article}
\usepackage[margin=1in]{geometry}
\usepackage{amsmath,amssymb}
\usepackage{booktabs}
\usepackage{multirow}
\usepackage{graphicx}
\usepackage{hyperref}
\usepackage{authblk}
\usepackage[numbers,sort]{natbib}

\title{ZoneBoost: A Fully Transparent, Zone-Based Alternative to Tree-Based Gradient Boosting}

\author[1]{Author One\thanks{Corresponding author.}}
\author[2]{Author Two}
\affil[1]{[Affiliation to be filled in]}
\affil[2]{[Affiliation to be filled in]}

\date{}

\begin{document}
\maketitle

\begin{abstract}
[Bilingual abstract delivered separately per project convention -- see accompanying abstract document.]
\end{abstract}

\section{Introduction}
\label{sec:intro}

Gradient boosting machines (GBMs) such as XGBoost \citep{chen2016xgboost}, LightGBM
\citep{ke2017lightgbm}, and CatBoost \citep{prokhorenkova2018catboost} are the
dominant tool for tabular prediction problems, routinely outperforming deep
learning on structured data. Their weak learner, however, is a decision tree:
a prediction is the sum of hundreds of tree traversals, each contributing a
leaf value shaped by an internal split sequence that is not directly
inspectable as a quantity a domain expert can audit line by line. The
standard response has been \emph{post hoc} explanation -- SHAP
\citep{lundberg2017shap} and LIME \citep{ribeiro2016lime} approximate, after
the fact, what a fitted black-box model is doing. Both are sampling-based
approximations: they estimate an attribution, they do not read one off the
model's own arithmetic.

A parallel line of work keeps the model itself additive and interpretable by
construction. Generalized Additive Models (GAMs) \citep{hastie1986gam} fit a
smooth, inspectable function per feature; the Explainable Boosting Machine
(EBM) \citep{lou2013ga2m,nori2019interpretml} boosts an ensemble of shallow
per-feature and per-feature-pair trees so that the resulting additive
function stays both accurate and glass-box. EBM is the closest existing model
to the one studied in this paper: both are additive, boosted, and expose
pairwise interactions as first-class, inspectable terms.

This paper studies \textbf{ZoneBoost}, a gradient boosting variant that
replaces the tree weak learner with a zone-grid weak learner built entirely
from descriptive statistics -- quantiles, group counts, and group averages.
Every boosting round splits each predictor's range into a small number of
data-driven zones (one zone per category for nominal features), computes an
empirical-Bayes-shrunk average residual per zone or zone pair, and combines
these zone-level statistics via Lasso stacking. No decision tree is ever fit;
no gradient-descent weight update happens in a hidden parameter space. A
prediction's exact decomposition into main effects and interaction terms is
available directly from the fitted object via an \texttt{explain()} method,
with no sampling and no approximation, unlike SHAP/LIME applied to a
tree-based GBM.

The central empirical question this paper asks is not whether ZoneBoost can
be made to win a leaderboard -- its value proposition is exact,
zero-approximation attribution, not necessarily state-of-the-art accuracy --
but rather: \emph{how large is the accuracy gap (if any) relative to
established GBMs and the closest interpretable competitor (EBM), which of
ZoneBoost's own architectural components actually drive whatever accuracy it
achieves, and does its interpretability/uncertainty-quantification machinery
do what it claims?} We answer this with:

\begin{itemize}
  \item A cross-validated benchmark against XGBoost, LightGBM, CatBoost, and
    EBM across six public datasets spanning regression, binary
    classification, multiclass classification, categorical-heavy features,
    and missing values (Section~\ref{sec:results}).
  \item A six-axis ablation study isolating the effect of empirical-Bayes
    shrinkage, zone-boundary treatment, interaction order, categorical
    encoding, row/column subsampling, and boundary smoothing
    (Section~\ref{sec:ablation}).
  \item Auxiliary experiments verifying explanation faithfulness against
    SHAP, conformal prediction interval coverage, classifier calibration, and
    drift detection on genuinely shifted data (Section~\ref{sec:auxiliary}).
\end{itemize}

\section{Related Work}
\label{sec:related}

\paragraph{Tree-based gradient boosting.} \citet{friedman2001greedy}
formalized gradient boosting as stagewise functional gradient descent.
XGBoost \citep{chen2016xgboost}, LightGBM \citep{ke2017lightgbm}, and CatBoost
\citep{prokhorenkova2018catboost} refined this with regularized objectives,
histogram-based splitting, and ordered boosting for categorical features,
respectively. All three remain the default choice for tabular prediction and
form this paper's ``raw accuracy'' comparison set.

\paragraph{Interpretable additive models.} GAMs \citep{hastie1986gam} fit
$g(y) = \sum_i f_i(x_i)$ with each $f_i$ a smooth univariate function.
\citet{lou2013ga2m} extended this to GA2M, adding a bounded set of pairwise
interaction terms selected by explained-variance gain -- the architecture EBM
\citep{nori2019interpretml} implements as a boosted ensemble of shallow
trees per term. ZoneBoost occupies the same additive-plus-pairwise-
interaction family but replaces EBM's per-term shallow-tree/spline fit with
zone-grid descriptive statistics under explicit empirical-Bayes shrinkage,
and replaces one-vs-rest multiclass boosting with a native multinomial
softmax booster.

\paragraph{Post-hoc explanation.} SHAP \citep{lundberg2017shap} and LIME
\citep{ribeiro2016lime} are model-agnostic, sampling-based attribution
methods applicable to any fitted model, including opaque GBMs. Because they
are approximations, their own faithfulness is not directly checkable against
a black-box model's true internal computation. ZoneBoost's \texttt{explain()}
sidesteps this by construction: the attribution \emph{is} the computation,
so Section~\ref{sec:auxiliary} can measure SHAP's approximation error against
a known-exact ground truth rather than merely reporting SHAP values on
faith.

\paragraph{Conformal prediction.} Split-conformal prediction gives
distribution-free coverage guarantees; Conformalized Quantile Regression
(CQR) \citep{romano2019cqr} combines this with quantile regression so
interval width varies with $X$, using the rearrangement fix of
\citet{chernozhukov2010rearrangement} to resolve quantile crossing.
ZoneBoost's \texttt{ConformalizedQuantileRegressor} is a direct application of
CQR built on two quantile-mode ZoneBoost fits.

\paragraph{Empirical-Bayes shrinkage.} ZoneBoost shrinks each zone's raw
mean toward a hierarchical prior via an m-estimate, with an optional
data-driven strength estimated by a DerSimonian-Laird-style
method-of-moments estimator \citep{dersimonian1986metaanalysis}, originally
developed for random-effects meta-analysis and reused here for the same
``compound normal means'' problem: pooling many noisy, unequally-sized group
estimates toward a common prior.

\section{Method: ZoneBoost Architecture}
\label{sec:method}

ZoneBoost fits an additive ensemble
$\hat{y}(x) = b + \sum_{r=1}^{R} \eta \cdot h_r(x)$, where $b$ is a baseline
statistic (the training mean, for squared-error loss), $\eta$ is the
learning rate, and each weak learner $h_r$ is fit on the current residual
using only zone-grid descriptive statistics -- never a tree, never a
gradient-descent weight update.

\paragraph{Zones.} For a continuous predictor, each round searches for the
cut point that most reduces within-zone residual variance, exactly the
criterion a regression tree stump uses, re-derived fresh every round from
that round's own residual rather than fixed once from quantile bins. For a
declared or auto-detected categorical predictor, every distinct value is its
own zone -- no cut-point search is performed, since there is no meaningful
notion of ``adjacent'' for a nominal label.

\paragraph{Interactions.} In addition to per-feature main effects, every
pair of subsampled columns is screened for a joint zone-grid interaction;
candidates are ranked by a cheap cross-fitted proxy and the strongest pairs
are fit exactly. A bounded, adaptive search for third-order interactions is
available (opt-in) but off by default, so the default configuration matches
every prior release's behavior exactly.

\paragraph{Empirical-Bayes shrinkage.} Each zone's raw cell mean
$\bar{y}_j$ (with $n_j$ supporting rows) is shrunk toward a prior $\mu_0$
(the global mean for a main effect, the already-shrunk lower-order marginal
sum for an interaction) via
$\hat{\mu}_j = (n_j \bar{y}_j + m \mu_0) / (n_j + m)$, where $m$
(\texttt{shrinkage\_m}) controls how many rows a zone needs before it is
trusted as much as its prior. $m$ can be fixed or estimated per round via a
DerSimonian-Laird-style method-of-moments procedure.

\paragraph{Combination and cross-fitting.} A round's terms are combined via
Lasso rather than an equal-weight average, so an irrelevant term is zeroed
by the L1 penalty instead of diluting a shared $1/n_\text{terms}$ share.
Each round's own training rows are scored with cross-fitted, out-of-fold
zone tables to avoid the same row's residual leaking into the statistic it
is then judged against.

\paragraph{Classification.} \texttt{ZoneBoostClassifier} reuses the
identical weak learner, boosting in log-odds space for binary targets
(a single sigmoid booster) and in a native multinomial softmax space for
$K \geq 3$ classes -- one booster jointly maintaining all $K$ logits per
row, rather than $K$ independent one-vs-rest boosters that never observe
each other's current score.

\paragraph{Uncertainty quantification.} \texttt{predict\_interval} gives a
single split-conformal margin; \texttt{ConformalizedQuantileRegressor} wraps
two quantile-mode fits (at $\alpha/2$ and $1-\alpha/2$) with CQR
calibration \citep{romano2019cqr} for a locally-adaptive interval whose
width varies with $X$.

\paragraph{Drift detection.} \texttt{compare\_models} is a stateless
comparison between two already-fitted regressors (e.g. last quarter's model
and this quarter's), reporting feature-importance drift, zone-boundary/
population shift, and prediction-shift statistics on a shared evaluation
set.

\section{Experimental Setup}
\label{sec:setup}

\paragraph{Datasets.} Six public datasets were used, chosen to span
regression and classification, continuous-only and categorical-heavy
feature sets, complete and missing data, and binary and multiclass targets
(Table~\ref{tab:datasets}). California Housing and Bike Sharing Demand were
subsampled to 4{,}000 rows (fixed seed) purely for cross-validated
wall-clock time across five models and six ablation axes -- a disclosed
speed tradeoff applied identically to every model, not a hyperparameter
change favoring any one of them, matching the methodology already used in
this project's own \texttt{benchmarks/compare\_lightgbm.py}.

\begin{table}[t]
\centering
\caption{Datasets used in the benchmark and ablation study.}
\label{tab:datasets}
\begin{tabular}{lllrrl}
\toprule
Dataset & Task & Source & $n$ & $p$ & Notes \\
\midrule
California Housing & Regression & scikit-learn & 4{,}000 & 8 & continuous only \\
Diabetes & Regression & scikit-learn & 442 & 10 & small $n$, continuous \\
Bike Sharing Demand & Regression & OpenML & 4{,}000 & 12 & categorical + continuous \\
Breast Cancer Wisconsin & Binary classification & scikit-learn & 569 & 30 & continuous only \\
Wine & Multiclass (3) classification & scikit-learn & 178 & 13 & tests native multinomial \\
Titanic & Binary classification & OpenML & 1{,}309 & 7 & categorical + missing values \\
\bottomrule
\end{tabular}
\end{table}

\paragraph{Models.} ZoneBoost is compared against XGBoost
\citep{chen2016xgboost}, LightGBM \citep{ke2017lightgbm}, CatBoost
\citep{prokhorenkova2018catboost}, and EBM \citep{nori2019interpretml}. Every
model is run at its own library's out-of-the-box defaults with only
\texttt{random\_state} fixed -- no hyperparameter tuning that favors any one
model, matching this project's existing disclosed benchmark philosophy. EBM
is run with \texttt{n\_jobs=1}: its default parallel backend spawns one
worker process per outer bag, and the fixed process-spawn overhead on the
evaluation machine dominated wall-clock time independent of real compute
cost -- an environment/parallelism-backend control, not a change to
\texttt{outer\_bags} or any other modeling default.

\paragraph{Protocol.} All results are 3-fold (regression: \texttt{KFold};
classification: \texttt{StratifiedKFold}) cross-validated, not a single
train/test split. Regression is scored by RMSE and $R^2$; classification by
accuracy, macro-F1, ROC AUC, and log loss.

\section{Results: Main Benchmark}
\label{sec:results}

Table~\ref{tab:regression} reports all three regression datasets;
Table~\ref{tab:classification} reports all three classification datasets.

\begin{table}[t]
\centering
\caption{Regression results (3-fold CV, mean $\pm$ std). Best RMSE per
dataset in bold.}
\label{tab:regression}
\begin{tabular}{llrrr}
\toprule
Dataset & Model & RMSE & $R^2$ & Fit time (s) \\
\midrule
\multirow{5}{*}{California Housing}
 & ZoneBoost & 0.5328 $\pm$ 0.0134 & 0.7836 $\pm$ 0.0037 & 7.90 \\
 & XGBoost   & 0.5482 $\pm$ 0.0096 & 0.7708 $\pm$ 0.0051 & 0.11 \\
 & LightGBM  & 0.5143 $\pm$ 0.0131 & 0.7984 $\pm$ 0.0044 & 0.69 \\
 & CatBoost  & \textbf{0.4949 $\pm$ 0.0176} & \textbf{0.8133 $\pm$ 0.0079} & 4.48 \\
 & EBM       & 0.5171 $\pm$ 0.0165 & 0.7962 $\pm$ 0.0064 & 197.49 \\
\midrule
\multirow{5}{*}{Diabetes}
 & ZoneBoost & 57.16 $\pm$ 2.00 & 0.4420 $\pm$ 0.0797 & 7.27 \\
 & XGBoost   & 65.00 $\pm$ 3.93 & 0.2832 $\pm$ 0.0731 & 0.08 \\
 & LightGBM  & 60.17 $\pm$ 4.44 & 0.3816 $\pm$ 0.1019 & 0.03 \\
 & CatBoost  & 58.96 $\pm$ 3.41 & 0.4091 $\pm$ 0.0687 & 3.70 \\
 & EBM       & \textbf{55.21 $\pm$ 1.65} & \textbf{0.4810 $\pm$ 0.0593} & 130.13 \\
\midrule
\multirow{5}{*}{Bike Sharing Demand}
 & ZoneBoost & 56.92 $\pm$ 1.61 & 0.9015 $\pm$ 0.0048 & 23.43 \\
 & XGBoost   & 50.38 $\pm$ 1.10 & 0.9228 $\pm$ 0.0045 & 0.20 \\
 & LightGBM  & 49.02 $\pm$ 1.21 & 0.9269 $\pm$ 0.0037 & 1.33 \\
 & CatBoost  & \textbf{47.12 $\pm$ 1.28} & \textbf{0.9325 $\pm$ 0.0032} & 75.95 \\
 & EBM       & 57.72 $\pm$ 0.99 & 0.8987 $\pm$ 0.0027 & 386.17 \\
\bottomrule
\end{tabular}
\end{table}

\begin{table}[t]
\centering
\caption{Classification results (3-fold stratified CV, mean $\pm$ std).
Best accuracy per dataset in bold.}
\label{tab:classification}
\begin{tabular}{llrrrrr}
\toprule
Dataset & Model & Accuracy & Macro-F1 & AUC & Log loss & Fit time (s) \\
\midrule
\multirow{5}{*}{Breast Cancer Wisconsin}
 & ZoneBoost & 0.9631 $\pm$ 0.0001 & 0.9602 & 0.9953 $\pm$ 0.0028 & 0.1196 & 52.15 \\
 & XGBoost   & \textbf{0.9701 $\pm$ 0.0163} & 0.9682 & 0.9924 $\pm$ 0.0040 & 0.1052 & 0.09 \\
 & LightGBM  & 0.9631 $\pm$ 0.0075 & 0.9604 & 0.9933 $\pm$ 0.0042 & 0.1454 & 1.08 \\
 & CatBoost  & 0.9684 $\pm$ 0.0113 & 0.9662 & 0.9931 $\pm$ 0.0048 & 0.0947 & 3.88 \\
 & EBM       & \multicolumn{5}{c}{not evaluated -- see Section~\ref{sec:limitations}} \\
\midrule
\multirow{5}{*}{Wine}
 & ZoneBoost & \textbf{0.9831 $\pm$ 0.0138} & 0.9839 & 0.9980 $\pm$ 0.0025 & 0.1055 & 27.27 \\
 & XGBoost   & 0.9496 $\pm$ 0.0362 & 0.9514 & 0.9967 $\pm$ 0.0025 & 0.1141 & 0.10 \\
 & LightGBM  & 0.9552 $\pm$ 0.0154 & 0.9571 & 0.9982 $\pm$ 0.0016 & 0.1287 & 1.23 \\
 & CatBoost  & 0.9440 $\pm$ 0.0205 & 0.9472 & 0.9959 $\pm$ 0.0029 & 0.1435 & 2.68 \\
 & EBM       & 0.9720 $\pm$ 0.0156 & 0.9731 & 0.9989 $\pm$ 0.0010 & 0.0838 & 863.64 \\
\midrule
\multirow{5}{*}{Titanic}
 & ZoneBoost & 0.7998 $\pm$ 0.0029 & 0.7816 & 0.8457 $\pm$ 0.0129 & 0.4543 & 4.30 \\
 & XGBoost   & 0.7945 $\pm$ 0.0159 & 0.7798 & 0.8441 $\pm$ 0.0060 & 0.5494 & 0.08 \\
 & LightGBM  & \textbf{0.8113 $\pm$ 0.0113} & 0.7979 & 0.8549 $\pm$ 0.0066 & 0.4932 & 0.94 \\
 & CatBoost  & 0.8060 $\pm$ 0.0088 & 0.7832 & \textbf{0.8557 $\pm$ 0.0103} & 0.4398 & 52.82 \\
 & EBM       & 0.7968 $\pm$ 0.0108 & 0.7799 & 0.8503 $\pm$ 0.0066 & 0.4515 & 21.32 \\
\bottomrule
\end{tabular}
\end{table}

On California Housing, a large-$n$, continuous-only dataset, ZoneBoost is
competitive with but does not beat the strongest tree-based competitor
(CatBoost, $R^2=0.813$ vs.\ ZoneBoost's $0.784$), while clearly beating
XGBoost and roughly matching EBM and LightGBM. On Diabetes -- a much smaller
dataset ($n=442$) where empirical-Bayes shrinkage toward a hierarchical
prior should matter most, since raw per-zone cell means are noisiest exactly
when a zone has few supporting rows -- ZoneBoost is the second-best model
overall ($R^2=0.442$), beating every tree-based competitor including
CatBoost, and trailing only EBM ($R^2=0.481$), the other model in this
comparison set built on an additive, shrinkage-aware architecture rather
than raw per-leaf tree statistics. This is consistent with the hypothesis
that ZoneBoost's advantage, where it has one, comes specifically from
small-sample regularization rather than raw representational capacity. On
Bike Sharing Demand, the one regression dataset with genuine categorical
features (season, weather, etc., each given their own zone rather than a
cut-point search), ZoneBoost trails all three tree-based competitors and
roughly matches EBM -- the categorical-heavy, large-$n$ regime is the one
place in the regression benchmark where ZoneBoost's zone-grid weak learner
does not close the gap to tree-based boosting.

On the classification side, ZoneBoost is essentially tied for best on
Breast Cancer Wisconsin (AUC $0.995$, a hair below XGBoost's best accuracy)
and is the outright best model on Wine by accuracy ($0.983$), a small
($n=178$), multiclass ($K=3$) dataset exercising ZoneBoost's native
multinomial softmax booster rather than a one-vs-rest reduction -- the
clearest win for ZoneBoost anywhere in the main benchmark. On Titanic, the
one classification dataset with both genuine categorical features and
missing values, ZoneBoost is mid-pack, trailing LightGBM and CatBoost on
accuracy and AUC but ahead of XGBoost and roughly matching EBM. Taken
together with the regression results, ZoneBoost's two clearest wins
(Diabetes, Wine) are both small-$n$ datasets, and its one clear loss (Bike
Sharing Demand) is the largest categorical-heavy dataset -- consistent with
the shrinkage-driven story developed further in
Section~\ref{sec:discussion}, though Titanic (categorical, missing values,
but only mid-pack rather than a clear loss) is a partial counterexample
worth flagging rather than smoothing over.

\section{Ablation Study}
\label{sec:ablation}

Six real, publicly-documented ZoneBoost parameters were toggled, one axis
at a time, holding every other setting at its default:

\begin{description}
  \item[A1 -- Empirical-Bayes shrinkage] \texttt{shrinkage\_m}$=0.01$
    (near-zero shrinkage; exactly $0$ triggers a $0/0$ empty-zone edge case
    in the current release, see Section~\ref{sec:limitations}) vs.\ the
    default $10$ vs.\ \texttt{learn\_shrinkage\_m=True} (data-driven
    strength, regressor only).
  \item[A2 -- Zone boundary treatment] the default adaptive per-round split
    search vs.\ boundaries fixed once via quantile binning (continuous
    columns pre-binned into 7 quantile bins and passed through the
    categorical zone path, so the same bin edges are reused every round
    instead of being re-derived).
  \item[A3 -- Interaction order] main-effects-only
    (\texttt{max\_pair\_interactions=0}) vs.\ the default pairwise
    (\texttt{max\_interaction\_order=2}) vs.\ pairwise-plus-triples
    (\texttt{max\_interaction\_order=3}).
  \item[A4 -- Categorical treatment] declared categorical, one zone per
    value (default) vs.\ ordinal-encoded and left continuous, subject to
    the adaptive cut-point search (only meaningful on Bike Sharing and
    Titanic, the two datasets with genuine categorical columns).
  \item[A5 -- Subsampling] the default (\texttt{row\_subsample=0.7},
    \texttt{col\_subsample=0.7}) vs.\ no subsampling (both $1.0$).
  \item[A6 -- Boundary smoothing]
    \texttt{adaptive\_boundary\_smoothing=False} (default, hard zone
    lookup) vs.\ \texttt{True} (cross-fitted soft interpolation).
\end{description}

Table~\ref{tab:ablation-reg} reports the regression ablation axes (all
three regression datasets, all six axes); the classification ablation axes
are reported separately below in Table~\ref{tab:ablation-clf}.

\begin{table}[t]
\centering
\small
\caption{Regression ablation study: RMSE per axis/variant/dataset (3-fold
CV). Default configuration in each axis marked with $^\dagger$.}
\label{tab:ablation-reg}
\begin{tabular}{llrrr}
\toprule
Axis & Variant & Cal.\ Housing & Diabetes & Bike Sharing \\
\midrule
\multirow{3}{*}{A1 shrinkage}
 & near-zero ($m=0.01$)   & 0.5370 & 57.13 & 56.84 \\
 & default ($m=10$)$^\dagger$ & 0.5328 & 57.16 & 56.92 \\
 & learned $m$             & 0.5430 & 57.18 & 58.41 \\
\midrule
\multirow{2}{*}{A2 zone boundary}
 & adaptive (per-round)$^\dagger$ & 0.5328 & 57.16 & 56.92 \\
 & fixed quantile bins      & 0.6433 & 59.87 & 88.64 \\
\midrule
\multirow{3}{*}{A3 interaction order}
 & main effects only        & 0.6056 & 56.79 & 100.84 \\
 & pairwise$^\dagger$        & 0.5328 & 57.16 & 56.92 \\
 & pairwise+triples          & 0.5312 & FAIL$^*$ & FAIL$^*$ \\
\midrule
A4 categorical treatment & own-zone$^\dagger$ / ordinal & --- & --- & 56.92 / 56.91 \\
\midrule
\multirow{2}{*}{A5 subsampling}
 & default (0.7/0.7)$^\dagger$ & 0.5328 & 57.16 & 56.92 \\
 & no subsampling (1.0/1.0)   & 0.5403 & 59.43 & 57.07 \\
\midrule
\multirow{2}{*}{A6 boundary smoothing}
 & hard lookup$^\dagger$      & 0.5328 & 57.16 & 56.92 \\
 & adaptive smoothing         & 0.5337 & 58.70 & 57.06 \\
\bottomrule
\end{tabular}

\vspace{2pt}
\footnotesize $^*$\texttt{max\_interaction\_order=3} raises an
\texttt{IndexError} on Diabetes and Bike Sharing Demand in the evaluated
release (v0.38.0); see Section~\ref{sec:limitations}. A4 is reported only
for Bike Sharing Demand, the one regression dataset with genuine
categorical columns.
\end{table}

By far the largest single-axis effect in the regression ablation is A2
(zone boundary treatment): freezing zone boundaries to fixed quantile bins
instead of re-deriving the split per round costs 21\% relative RMSE on
California Housing (0.533 $\to$ 0.643), 5\% on Diabetes, and a striking
56\% on Bike Sharing Demand (56.9 $\to$ 88.6) -- the adaptive, per-round
split search is not a minor refinement but the single largest contributor
to ZoneBoost's accuracy of any axis tested. A3 (interaction order) shows
the expected direction -- main-effects-only is markedly worse than pairwise
on California Housing and (especially) Bike Sharing Demand -- confirming
that the pairwise interaction terms are doing real work, not just adding
inspectable-but-inert structure. A1 (shrinkage) and A5 (subsampling) show
smaller, dataset-dependent effects: near-zero shrinkage is essentially free
on Diabetes and Bike Sharing Demand but costs a little on California
Housing, while disabling subsampling entirely (row/col at 1.0) modestly
\emph{hurts} both California Housing and Diabetes -- subsampling here
functions as intended regularization, not merely a speed optimization. A6
(boundary smoothing) has the smallest effect of any axis tested, and
learned shrinkage strength (A1's third arm) underperforms the fixed
default on all three regression datasets, suggesting the method-of-moments
estimator is not (yet) pulling its weight relative to just using a fixed
$m=10$.

\begin{table}[t]
\centering
\small
\caption{Classification ablation study: accuracy per axis/variant/dataset
(3-fold stratified CV). Default configuration in each axis marked with
$^\dagger$.}
\label{tab:ablation-clf}
\begin{tabular}{llrrr}
\toprule
Axis & Variant & Breast Cancer W. & Wine & Titanic \\
\midrule
\multirow{2}{*}{A1 shrinkage}
 & near-zero ($m=0.01$)     & 0.9613 & 0.9775 & 0.8021 \\
 & default ($m=10$)$^\dagger$ & 0.9631 & 0.9831 & 0.7998 \\
\midrule
\multirow{2}{*}{A2 zone boundary}
 & adaptive (per-round)$^\dagger$ & 0.9631 & 0.9831 & 0.7998 \\
 & fixed quantile bins       & 0.9596 & 0.9776 & 0.8029 \\
\midrule
\multirow{3}{*}{A3 interaction order}
 & main effects only         & 0.9684 & 0.9720 & 0.7846 \\
 & pairwise$^\dagger$         & 0.9631 & 0.9831 & 0.7998 \\
 & pairwise+triples           & FAIL$^*$ & 0.9831 & FAIL$^*$ \\
\midrule
A4 categorical treatment & own-zone$^\dagger$ / ordinal & --- & --- & 0.7998 / 0.7991 \\
\midrule
\multirow{2}{*}{A5 subsampling}
 & default (0.7/0.7)$^\dagger$ & 0.9631 & 0.9831 & 0.7998 \\
 & no subsampling (1.0/1.0)    & 0.9613 & 0.9719 & 0.8014 \\
\midrule
\multirow{2}{*}{A6 boundary smoothing}
 & hard lookup$^\dagger$       & 0.9631 & 0.9831 & 0.7998 \\
 & adaptive smoothing          & 0.9648 & 0.9775 & 0.8029 \\
\bottomrule
\end{tabular}

\vspace{2pt}
\footnotesize $^*$\texttt{max\_interaction\_order=3} raises an
\texttt{IndexError} on Breast Cancer Wisconsin and Titanic in the
evaluated release (v0.38.0), the same failure mode documented for the
regression axis; Wine (the one classification dataset where this arm ran
successfully) shows no measurable change from the pairwise default. A4 is
reported only for Titanic, the one classification dataset with genuine
categorical columns.
\end{table}

The classification ablation broadly confirms the regression ablation's
central finding: A2 (adaptive zone boundaries) and A5 (subsampling) are
again real, measurable regularizers rather than incidental defaults --
disabling either costs accuracy on Breast Cancer Wisconsin and Wine, the
two datasets with the largest per-class sample sizes to regularize over.
Titanic is the one dataset where several axes point the \emph{other}
direction (fixed quantile bins, no subsampling, and adaptive smoothing all
very slightly \emph{improve} accuracy over the default), though every one
of these differences is a fraction of a percentage point and well within
the fold-to-fold noise visible in Table~\ref{tab:classification}'s standard
deviations -- consistent with Titanic being the smallest, noisiest
classification dataset ($n=1{,}309$ but with a 30\% missing-\texttt{age}
column and a famously high-variance target) rather than evidence that
regularization actively hurts there. A3 (interaction order) is largely
uninformative on the classification side because the triples arm fails to
evaluate on two of the three datasets; on Wine, the one dataset where it
did run, pairwise-plus-triples is statistically indistinguishable from
pairwise alone, offering no evidence that third-order interactions would
help even where they can be measured. A4 (categorical treatment on
Titanic) and A1 (shrinkage) both show effects under half a percentage
point in either direction -- the smallest-effect axes on the classification
side, mirroring A6 being the smallest-effect axis on the regression side.

\section{Auxiliary Experiments}
\label{sec:auxiliary}

\subsection{Explanation faithfulness}
\texttt{explain()}'s per-term contributions are checked to sum exactly to
\texttt{predict()}'s output (to float tolerance) on held-out rows of
California Housing and Breast Cancer Wisconsin, then aggregated to
per-raw-feature attributions (an interaction term's contribution split
evenly between its two constituent features) and compared against a
SHAP \texttt{KernelExplainer} approximation of the \emph{same fitted model}
-- a faithfulness check on SHAP's own approximation quality against a known-
exact ground truth, not merely a report of SHAP values on faith.

On held-out California Housing rows, \texttt{explain()}'s per-term
contributions sum to \texttt{predict()}'s output to within
$7.1\times10^{-15}$ absolute error -- floating-point precision, i.e.\
exact, as designed. The same holds in log-odds space on Breast Cancer
Wisconsin ($4.4\times10^{-15}$). A SHAP \texttt{KernelExplainer}
approximation of the \emph{same fitted} California Housing model, evaluated
on 30 held-out rows (200 coalition samples per row), achieves a mean
per-row Pearson correlation of only $0.868$ against ZoneBoost's exact
per-raw-feature attribution, with a mean L1 attribution gap of $0.676$
(target units) per row -- a real, non-trivial approximation error, not
merely noise, on a case where the ground truth is actually known and
checkable rather than assumed. This is the paper's central interpretability
finding: SHAP's own approximation quality, typically taken on faith because
no ground truth exists for the black-box models it is normally applied to,
is here directly measurable and imperfect.

\subsection{Conformal interval coverage}
\texttt{ConformalizedQuantileRegressor} at $\alpha=0.1$ (target 90\%
coverage) is evaluated on held-out data from all three regression datasets,
reporting empirical coverage and interval width (mean and standard
deviation -- a nonzero width standard deviation is direct evidence of local
adaptivity, as opposed to a single constant-width margin).

\begin{table}[h]
\centering
\caption{Conformalized quantile regression: empirical coverage at $\alpha=0.1$ (target 90\%).}
\label{tab:conformal}
\begin{tabular}{lrrr}
\toprule
Dataset & Empirical coverage & Mean width & Width std.\ \\
\midrule
California Housing & 0.890 & 1.871 & 0.597 \\
Diabetes & 0.895 & 184.11 & 25.50 \\
Bike Sharing Demand & 0.881 & 213.58 & 122.74 \\
\bottomrule
\end{tabular}
\end{table}

All three datasets land close to, and consistently slightly below, the
90\% target (0.881--0.895) -- a small, uniform undercoverage rather than a
one-off, worth flagging for future work rather than dismissing as noise. In
every case the interval width's standard deviation is a substantial
fraction of its mean, confirming the interval is genuinely locally
adaptive (widest on Bike Sharing Demand, where per-hour demand variance
itself varies sharply with covariates) rather than a constant-width band
dressed up as one.

\subsection{Classifier calibration}
\texttt{ZoneBoostClassifier} is fit on Titanic with
\texttt{calibrate=False} and \texttt{calibrate=True}, reporting Brier score
and expected calibration error (ECE, 10 quantile bins) for both.

On Titanic, \texttt{calibrate=False} gives Brier $=0.1358$, ECE $=0.0692$;
\texttt{calibrate=True} gives Brier $=0.1392$, ECE $=0.0724$ -- both
slightly \emph{worse} than the uncalibrated model on this particular
dataset and split. This is a genuine, disclosed null result rather than
the expected direction: post-hoc isotonic calibration is not guaranteed to
help, particularly when the calibration split itself is small (Titanic,
$n=1{,}309$, a 30\% held-out test split) and the base classifier is already
reasonably well-calibrated by virtue of its own additive, shrinkage-based
architecture.

\subsection{Drift detection}
Two \texttt{ZoneBoostRegressor} models are fit on a real covariate-shift
split of California Housing (older vs.\ newer housing stock, split at the
median \texttt{HouseAge}, not a synthetic perturbation), and
\texttt{compare\_models} is used to report the largest feature-importance
shifts between them.

\texttt{compare\_models} correctly surfaces \texttt{AveOccup} (average
household occupancy) as the largest importance shift between the
older-housing and newer-housing models (0.063 $\to$ 0.260, a
$4.1\times$ increase), followed by \texttt{MedInc} (0.334 $\to$ 0.407) and
two interaction terms newly significant in the newer-housing model
(\texttt{AveOccup x Latitude}, \texttt{HouseAge x Longitude}). The newer
model also fits noticeably better on its own data (RMSE $0.500$ vs.\
$0.664$ for the older model evaluated out-of-sample on the same newer
split), consistent with a genuine, not merely cosmetic, distributional
difference between the two housing-age segments -- exactly the kind of
signal a drift-monitoring tool is meant to catch.

\section{Discussion}
\label{sec:discussion}

The main-benchmark pattern (Section~\ref{sec:results}) was initially
consistent with a specific, falsifiable story about \emph{where}
ZoneBoost's empirical-Bayes shrinkage earns its keep: on the smaller
dataset (Diabetes, $n=442$), where a zone's raw cell mean is estimated from
few rows and shrinkage toward a hierarchical prior meaningfully denoises
it, ZoneBoost beats every tree-based competitor and trails only EBM -- the
other shrinkage-aware additive model in the comparison set. On the larger,
continuous-only California Housing ($n=4{,}000$), where individual zones
are populated by enough rows that shrinkage has less raw material to
correct, ZoneBoost's advantage narrows to roughly matching EBM/LightGBM
and trailing CatBoost. The ablation study (Section~\ref{sec:ablation}) was
designed specifically to confirm or refute shrinkage as the mechanism
behind this pattern, rather than merely assert it -- and the ablation
\emph{refutes} it. Across every regression and classification dataset
tested, near-zero shrinkage (A1, $m=0.01$) performs statistically
indistinguishably from, and on four of the nine dataset/task combinations
slightly \emph{better} than, the default ($m=10$); if anything, Diabetes
shows one of the \emph{smallest} shrinkage effects of any dataset tested
(57.13 vs.\ 57.16 RMSE, near-zero vs.\ default), the opposite of what the
original hypothesis predicted. Empirical-Bayes shrinkage, as implemented
and at its default strength, is not the mechanism behind ZoneBoost's
small-$n$ competitiveness.

What the ablation identifies instead, with a much larger and more
consistent effect size across every dataset in both tasks, is A2: the
adaptive, per-round zone-boundary search. Freezing zone boundaries to
fixed quantile bins costs 5--56\% relative RMSE across the three regression
datasets and 0.6--3\% relative accuracy across the three classification
datasets, making it the single largest-effect axis found anywhere in this
study, regression or classification. This reframes the small-$n$ result:
Diabetes's win over every tree-based competitor is more plausibly explained
by ZoneBoost re-deriving its cut points fresh from each round's own
residual (rather than a shrinkage effect specific to small samples), since
this same mechanism is the dominant one on every dataset regardless of
$n$. The Diabetes/EBM comparison in particular is likely confounded by
other architectural differences (EBM's smooth per-feature splines vs.\
ZoneBoost's discrete zone-grid steps) rather than isolating a shrinkage
effect, since ZoneBoost's own ablation rules shrinkage out as the load-
bearing mechanism in its own architecture.

Synthesizing across all six datasets: ZoneBoost's two clearest positive
results are Diabetes (small-$n$ regression) and Wine (small-$n$,
multiclass classification exercising the native softmax booster), and its
one clear negative result is Bike Sharing Demand, the largest and most
categorical-heavy regression dataset, where the same zone-boundary
adaptivity that helps elsewhere is presumably outcompeted by tree-based
models' ability to place many more, deeper splits per round. Titanic
(categorical, missing values, smallest classification $n$) and California
Housing (large, continuous-only) both land mid-pack rather than at either
extreme. Practically, this suggests preferring ZoneBoost over an opaque
GBM specifically when exact, zero-approximation attribution is a hard
requirement \emph{and} the data is small-to-moderate $n$ or multiclass
with genuinely interacting features -- not as a general-purpose accuracy
play, and not (yet) for large, categorical-heavy regression, where the
measured gap to tree-based boosting is largest.

\section{Limitations}
\label{sec:limitations}

ZoneBoost's own documentation is explicit that its value proposition is
exact attribution, not necessarily state-of-the-art tabular accuracy; this
paper's benchmark is designed to measure that gap honestly rather than
assume it away. All datasets were subsampled or used at moderate scale
(largest $n=4{,}000$) for cross-validated wall-clock tractability across
five models and six ablation axes; results at larger $n$ (millions of rows,
the regime tree-based GBMs are most optimized for) are not evaluated here.
EBM was run with \texttt{n\_jobs=1} for environment reasons documented in
Section~\ref{sec:setup}; this affects wall-clock time only, not its fitted
predictions. The interaction-attribution split used for the SHAP comparison
(50/50 credit between an interaction's two features) is a disclosed
convention, not the only defensible one.

Two implementation edge cases in the evaluated release (v0.38.0) were
discovered during this study and are disclosed rather than silently
routed around. First, \texttt{shrinkage\_m=0} (exact zero, as opposed to
the near-zero $0.01$ used in Section~\ref{sec:ablation}'s A1 arm) produces
a $0/0$ division for any zone with no supporting rows in a given round,
raising an error rather than returning a prediction; the ablation study
therefore reports the near-zero arm instead of a true zero-shrinkage arm.
Second, \texttt{max\_interaction\_order=3} raises an \texttt{IndexError}
on four of the six datasets tested (Diabetes, Bike Sharing Demand, Breast
Cancer Wisconsin, and Titanic) but not the other two (California Housing
and Wine), with no sample-size threshold that cleanly separates the two
groups ($n$ ranges from 442 to 4{,}000 among the failures and from 178 to
4{,}000 among the successes) -- so this ablation arm is reported as not
evaluable on the four affected datasets rather than papered over with a
substitute number, and the maintainers should treat the failure condition
as still uncharacterized rather than simply ``small $n$.''
Both are reported here as findings for the maintainers, not worked around
in the shipped library for this paper.

Third, and unrelated to ZoneBoost itself: the EBM baseline's fit on Breast
Cancer Wisconsin ($n=569$, $p=30$) did not complete within a $>$3-hour
wall-clock budget on the evaluation machine and was abandoned, versus
$\leq$45 minutes for every other (dataset, EBM) pair in this study --
consistent in direction with, but far more severe than, the per-outer-bag
process-spawn overhead already documented in Section~\ref{sec:setup}, and
plausibly compounded by Breast Cancer Wisconsin's comparatively high
feature count ($p=30$ vs.\ $p\leq13$ for every other dataset EBM completed
on), which enlarges EBM's default pairwise-interaction candidate screen
roughly quadratically. This cell is reported in
Table~\ref{tab:classification} as not evaluated rather than substituted
with a partial or estimated number.

\section{Conclusion}
\label{sec:conclusion}

ZoneBoost replaces the tree weak learner in gradient boosting with a
zone-grid weak learner built entirely from quantiles, group counts, and
empirical-Bayes-shrunk group averages, giving every prediction an exact,
zero-approximation decomposition into main-effect and interaction terms.
Across the six datasets evaluated here, its accuracy trails tree-based
competitors most on large, categorical-heavy regression data (Bike Sharing
Demand) but is competitive with, and on two small-to-moderate datasets
(Diabetes, Wine) better than, every tree-based competitor tested. A
six-axis ablation study designed to isolate the mechanism behind this
pattern rules out the hypothesized explanation -- empirical-Bayes
shrinkage -- and instead identifies the adaptive, per-round zone-boundary
search as by far the largest-effect architectural component, degrading
RMSE by up to 56\% and accuracy by up to 3\% when disabled, more than any
other axis tested in either task. This is itself a useful result: it
redirects where future architectural effort on ZoneBoost should go, and
illustrates why an ablation study earns its place in an interpretable-model
paper even when (especially when) it refutes the paper's own working
hypothesis. Its exact attribution is not merely a documentation claim: a
SHAP approximation of the same fitted model recovers it only partially
(mean correlation $0.868$), directly demonstrating that post-hoc
explanation carries a real, measurable cost that ZoneBoost's architecture
avoids by construction. Its conformal, calibration, and drift-detection
utilities largely deliver on their stated guarantees, with one clear
exception (calibration did not improve Brier score or ECE on Titanic)
reported as a genuine null result. Taken together, these results position
ZoneBoost as a specific-purpose tool: prefer it over an opaque GBM when
exact, auditable attribution is a hard requirement and the data is
small-to-moderate $n$ or multiclass, and prefer tree-based boosting when
raw accuracy on large, categorical-heavy data is the only objective.

\section*{Data and Code Availability}
ZoneBoost is open source (MIT license) and published on PyPI
(\texttt{pip install zoneboost}); source is available at
\url{https://github.com/stainaz/zoneboost}. All datasets used are public
(scikit-learn bundled datasets and OpenML). Experiment scripts and raw
results for this paper are available in the same repository, under
\texttt{paper/} (\url{https://github.com/stainaz/zoneboost/tree/main/paper}).

\section*{Ethics Declaration}
This study uses only public, previously de-identified benchmark datasets
and involves no human subjects research beyond secondary analysis of
publicly available data (e.g. Titanic passenger records, a standard,
long-public benchmark dataset).

\section*{Author Contributions (CRediT)}
[To be completed by the authors.]

\section*{Conflict of Interest}
The authors declare no conflict of interest.

\section*{Funding}
This research received no external funding.

\section*{AI Usage Disclosure}
[To be completed per target venue policy -- see accompanying disclosure
statement.]

\bibliographystyle{IEEEtranN}
\bibliography{references}

\end{document}
